\documentclass{beamer}
\usepackage[spanish]{babel}
\usepackage[utf8]{inputenc}
\usepackage{graphicx}

\newtheorem{definicion}{Definición}
\newtheorem{ejemplo}{Ejemplo}

%%%%%%%%%%%%%%%%%%%%%%%%%%%%%%%%%%%%%%%%%%%%%%%%%%%%%%%%%%%%%%%%%%%%%%%%%%%%%%%
\title[presentacion]{Introduccion al \Latex}
\author[L. Diaz]{Lorena Díaz Morales}
\date[15-03-2013]{15 de marzo de 2013}
%%%%%%%%%%%%%%%%%%%%%%%%%%%%%%%%%%%%%%%%%%%%%%%%%%%%%%%%%%%%%%%%%%%%%%%%%%%%%%%

\usetheme{Madrid}
%\usetheme{Antibes}
%\usetheme{boxes}
%\usetheme{tree}
%\usetheme{classic}

%%%%%%%%%%%%%%%%%%%%%%%%%%%%%%%%%%%%%%%%%%%%%%%%%%%%%%%%%%%%%%%%%%%%%%%%%%%%%%%
\begin{document}
  
%++++++++++++++++++++++++++++++++++++++++++++++++++++++++++++++++++++++++++++++  
\begin{frame}

  \includegraphics[width=0.15\textwidth]{img/ullesc.eps}
  \hspace*{7.5cm}
  \includegraphics[width=0.16\textwidth]{img/fmatesc.eps}
  \titlepage

  \begin{scriptsize}
    \begin{center}
     Facultad de Matemáticas \\
     Universidad de La Laguna
    \end{center}
  \end{scriptsize}

\end{frame}
%++++++++++++++++++++++++++++++++++++++++++++++++++++++++++++++++++++++++++++++  

%++++++++++++++++++++++++++++++++++++++++++++++++++++++++++++++++++++++++++++++  
\begin{frame}
  \frametitle{Índice}  
  \tableofcontents[pausesections]
\end{frame}
%++++++++++++++++++++++++++++++++++++++++++++++++++++++++++++++++++++++++++++++  


\section{Primera Sección}


%++++++++++++++++++++++++++++++++++++++++++++++++++++++++++++++++++++++++++++++  
\begin{frame}

\frametitle{¿Que es \LaTeX{}?}

\begin{definicion}
\Latex es un sistema de composicion de textos, orientado especialmente a la creación de libros,documentos científicos y técnicos que 
contengan fórmulas matemáticas como por ejemplo 
\[ x+y^{3t+7}=\sqrt{(b^2-4ac)/2a} \] \[ S_n=\sum_{i=1} a_i \]

\end{definicion}

\end{frame}
%++++++++++++++++++++++++++++++++++++++++++++++++++++++++++++++++++++++++++++++  

%++++++++++++++++++++++++++++++++++++++++++++++++++++++++++++++++++++++++++++++  
\begin{frame}

\frametitle{Formulas matamaticas}

\begin{block}{Ejemplo}
  \begin{itemize}
  \item
  $3x+4x-3x^3=0$ 
  \pause

  \item
  $\int sen(x) dx$
  \pause

  \item
  $\int\int_D x^2$
  

  \end{itemize}
\end{block}

\end{frame}
%++++++++++++++++++++++++++++++++++++++++++++++++++++++++++++++++++++++++++++++  

\section{¿Por que utilizar Latex{}?}
\begin{frame}
\begin{itemize}
\item
PROS
Latex es ampliamente utilizado en entornos cientificos. Muchas revistas aceptan documentos escritos en \Latex
\pause

\item
Excelente calidad del documento final con salida en distintos formatos:DVI,PDF,...
\pause

\item
Los ficheros fuente .tex son ficheros ASCII y pueden ser compilados en cualquier sistema operativo

\pause

\item
Es gratuito.
\pause

\item
Muy potente.
\pause

CONTRAS

\item
No es un procesador del tipo WYSIWYG "lo que escribes es lo que consigues", por lo que es necesario un procesador de compilación(con posibles errores...)
\pause
\end{itemize}
\end{frame}

%++++++++++++++++++++++++++++++++++++++++++++++++++++++++++++++++++++++++++++++  
\begin{frame}
\frametitle{Personalizando los nombres de caption y bibliografia}

Podemos cambiar los nombres y la numeración de las caption en los entornos de figuras y tablas. una forma sencilla es utilizar el paquete caption.

\end{frame}
%++++++++++++++++++++++++++++++++++++++++++++++++++++++++++++++++++++++++++++++  


\begin{frame}
  \frametitle{Bibliografía}

  \begin{thebibliography}{10}

    \beamertemplatebookbibitems
    \bibitem[URL: CTAN]ĺatex}
    {\small $http://www.pcmap.unizar.es/~pilar/latex.pdf$}

    \beamertemplatebookbibitems
    \bibitem[URL: CTAN]{latex} 
    CTAN. {\small $http://es.wikipedia.org/wiki/latex$}

  \end{thebibliography}
\end{frame}


%++++++++++++++++++++++++++++++++++++++++++++++++++++++++++++++++++++++++++++++  

\end{document}









